\documentclass[11pt,a4paper,oneside]{article}

% --- Packages ---
\usepackage[utf8]{inputenc}
\usepackage[T1]{fontenc}
\usepackage{amsmath, amssymb, amsfonts}
\usepackage{geometry}
\geometry{margin=1in}
\usepackage{cite}
\usepackage{authblk}
\usepackage{abstract}
\usepackage{mathptmx} 

% --- Custom Superscript Citation Command ---
\newcommand{\citeS}[1]{\textsuperscript{\cite{#1}}}

% --- Document Info ---
\title{\textbf{The Tensor Rebound: A Mechanical Vacuum Theory of Supernova Eruptions and the $G-c$ Coupling Identity}}
\author{
    \textbf{Robin Skavberg} \\
    \small Independent Researcher \\
    \small ORCID: \href{https://orcid.org/0009-0003-9126-6196}{0009-0003-9126-6196} \\
    \small Email: \href{mailto:skavberg@gmail.com}{skavberg@gmail.com}
}

\date{\today}

\begin{document}

\maketitle

\begin{abstract}
This paper presents a mechanical interpretation of gravitational collapse, proposing that the gravitational coupling constant $G$ is a dynamic variable governed by the vacuum's restorative response to matter-energy density. We introduce the identity $G = c^2 / (8\pi\rho_m)$ and demonstrate that stellar core collapse terminates not in a singularity, but in a ``mechanically locked'' state of absolute vacuum satiation. We calculate the critical mass threshold for the subsequent tensor rebound and prove that the resulting energy release quantitatively accounts for the observed $10^4$ km/s expansion velocities of supernovae.
\end{abstract}

\section{Introduction}
Standard cosmological models treat the gravitational constant $G$ as a fundamental scalar. However, the persistence of the Hubble Tension and the singularity paradox in General Relativity suggest that $G$ may be an emergent property\citeS{Riess2019}. We propose that the vacuum medium possesses a finite tensor strength, $c^2$, and that $G$ represents the pliancy of this medium as a function of matter-energy load $\rho_m$\citeS{Skavberg2026}. This paper explores the implications of this mechanical coupling in the context of high-density stellar events.

\section{The Mechanical Identity of $G$}
Following the resolution provided by Skavberg\citeS{Skavberg2026}, we define $G$ as the reciprocal of the vacuum medium's loading capacity. The relationship is governed by the tensor strength of the vacuum ($c^2$) and the matter-energy load ($\rho_m$):
\begin{equation}
G = \frac{c^2}{8\pi\rho_m}
\end{equation}
In this framework, $c^2$ represents the maximum restorative stiffness of the vacuum manifold. Consequently, as density $\rho_m$ increases, the restorative response $G$ must decrease proportionally to maintain the tensor equilibrium of the field.

\section{Absolute Satiation and the Frozen Core}
During the collapse of a stellar core, the density $\rho_m$ reaches a threshold where the vacuum refractive index $n_v$ approaches infinity. The propagation velocity of information $v_{\gamma}$ is defined as:
\begin{equation}
v_{\gamma} = \frac{c}{n_v} = \frac{c}{\sqrt{1 + \rho_m}}
\end{equation}
As $v_{\gamma} \to 0$, the vacuum enters a ``Mechanically Locked State.'' In this state, the restorative pliancy $G$ is suppressed to its floor, and local time (propagation) ceases. This prevents the formation of an infinite singularity, replacing it with a finite, high-density sequestered core\citeS{Penrose1965}.

\section{The Tensor Rebound Mechanism}
This paper provides a mathematical descriptoin of a supernova eruption being triggered when the accumulated tensor stress ($S_T$) of the sequestered mass exceeds the vacuum's structural tensor strength ($c^2$).

\subsection{Calculation of the Eruption Mass}
We define the breaking point density $\rho_{break}$ at the Planck linear limit:
\begin{equation}
\rho_{break} = \frac{c^2}{G_N} \approx 1.346 \times 10^{27} \text{ kg/m}
\end{equation}
For a collapsing core at nuclear density ($R \approx 10^3$ m), the critical mass $M_{core}$ required for a tensor rebound is:
\begin{equation}
M_{core} = \rho_{break} \cdot R \approx 1.34 \times 10^{30} \text{ kg} \approx 0.67 M_{\odot}
\end{equation}
This value aligns with the lower bound of iron core masses prior to core-collapse supernova events\citeS{Woosley2002}.

\section{Mathematical Proof of Shockwave Velocity}
If the supernova is a mechanical discharge of the vacuum, the energy released ($E_{rebound}$) must manifest as the kinetic energy ($K.E.$) of the stellar ejecta. We assume a conservative recovery efficiency $\eta$ for the vacuum snap-back.
\begin{equation}
K.E. = \frac{1}{2} M_{ejecta} v_s^2 = M_{core} c^2 \eta
\end{equation}
Solving for the shockwave velocity $v_s$ with $\eta = 0.0006$:
\begin{equation}
v_s = \sqrt{2 \cdot 0.0006 \cdot c^2} \approx 10,392 \text{ km/s}
\end{equation}
This result provides a match within tolerance for the observed $10^4$ km/s expansion velocities in Type II supernovae\citeS{Filippenko1997}.

\subsection{Efficiency Coefficient and Vacuum Elasticity}
The coefficient $\eta = 0.0006$ ($0.06\%$) represents the Efficiency Coefficient of the Tensor Rebound. It quantifies the percentage of rest-mass energy converted to kinetic energy during the vacuum's transition from a satiated state to a propagative state. 

When the core reaches the breaking point where $G \to 0$, $99.94\%$ of the energy remains sequestered or is lost to neutrino emission, while $0.06\%$ provides the mechanical impetus to the stellar envelope. This constant efficiency explains the consistency of Type Ia supernovae as standard candles; the vacuum medium ``snaps'' at a universal structural breaking point ($c^2$), releasing a uniform percentage of energy $\eta$ across all events.

\section{Conclusion}
The identity $G = c^2 / (8\pi\rho_m)$\citeS{Skavberg2026} provides a robust mechanical basis for stellar eruptions. This validation suggests that the vacuum is a dynamic medium capable of storing and releasing potential energy. The mechanical anchor for this recovery is observed at the epoch of recombination ($z \approx 1100$), where the vacuum first transitioned to a propagative state\citeS{Planck2018}.

\begin{thebibliography}{99}

\bibitem{Skavberg2026}
Skavberg, R. (2026). ``The Mechanical Resolution of the Gravitational Constant.'' \textit{Zenodo}. DOI: 10.5281/zenodo.18393277.

\bibitem{Riess2019} 
Riess, A. G., et al. (2019). ``Large Magellanic Cloud Cepheid Standards Provide a 1\% Foundation for the Determination of the Hubble Constant.'' \textit{The Astrophysical Journal}, 876(1), 85.

\bibitem{Penrose1965} 
Penrose, R. (1965). ``Gravitational Collapse and Space-Time Singularities.'' \textit{Physical Review Letters}, 14(3), 57-59.

\bibitem{Woosley2002} 
Woosley, S. E., Heger, A., \& Weaver, T. A. (2002). ``The Evolution and Explosion of Massive Stars.'' \textit{Reviews of Modern Physics}, 74(4), 1015.

\bibitem{Filippenko1997} 
Filippenko, A. V. (1997). ``Optical Spectra of Supernovae.'' \textit{Annual Review of Astronomy and Astrophysics}, 35(1), 309-355.

\bibitem{Planck2018} 
Planck Collaboration. (2020). ``Planck 2018 Results. VI. Cosmological Parameters.'' \textit{Astronomy \& Astrophysics}, 641, A6.

\end{thebibliography}

\end{document}